%
% Halaman Abstrak
%
% @author  Andreas Febrian @version 2.2.0 @edit by Ichlasul Affan
%

\chapter*{Abstrak}
\singlespacing

\noindent \begin{tabular}{l l p{10cm}} \ifx\blank\npmDua Nama&: & \penulisSatu
	\\
		Program Studi&: & \programSatu \\
	\else
		Nama Penulis 1 / Program Studi&: & \penulisSatu~/ \programSatu\\
		Nama Penulis 2 / Program Studi&: & \penulisDua~/ \programDua\\
	\fi
	\ifx\blank\npmTiga\else Nama Penulis 3 / Program Studi&: & \penulisTiga~/
		\programTiga\\
	\fi
	Judul&: & \judul \\
	Pembimbing&: & \pembimbingSatu \\
	\ifx\blank\pembimbingDua
    \else
        \ &\ & \pembimbingDua \\
    \fi
    \ifx\blank\pembimbingTiga
    \else
    	\ &\ & \pembimbingTiga \\
    \fi
\end{tabular} \\

\vspace*{0.5cm}

\noindent Kerja praktik ini dilaksanakan di PT TASPEN (Persero) dengan lingkup
pekerjaan pengembangan aplikasi internal Easy Mobile yang berfungsi menunjang
aktivitas karyawan, seperti pengajuan cuti dan Surat Perintah Tugas (SPT).
Selama kerja praktik berlangsung, pelaksana berperan sebagai Mobile Developer
Intern di Divisi Teknologi Informasi dengan tanggung jawab mengembangkan fitur
Request SPT dan Request Cuti menggunakan Flutter, serta memanfaatkan Figma,
GitLab, Postman, dan Android Studio sebagai pendukung. Proses pelaksanaan
menghadapi tantangan berupa keterbatasan dokumentasi API dan kerumitan user
flow, tetapi dapat diselesaikan melalui komunikasi intensif dengan mentor,
analisis mandiri, dan mengadopsi praktik clean architecture. Kegiatan ini tidak
hanya memperkuat kemampuan teknis dalam pengembangan aplikasi mobile, tetapi
juga mengasah keterampilan nonteknis seperti komunikasi, kerja sama tim,
adaptasi, serta pengelolaan waktu yang menjadi bekal penting dalam dunia kerja profesional. \\

\vspace*{0.2cm}

\noindent Kata kunci: \\ \f{mobile development}, Flutter, \f{front-end},
aplikasi perusahaan \\

\setstretch{1.4}
\newpage
